\section{Stopped routes}\label{sec:obstructions}

We group the stopped routes by the kind of stop. Entries labeled Research record summarize archived arguments whose full definitions and proofs are in the claim-index evidence; those summaries are not self-contained exclusion theorems. Elementary counterexamples below are stated directly.
\begin{itemize}[leftmargin=*]
\item A \emph{proved obstruction} is a theorem about an explicitly delimited class of constructions. Nothing is claimed outside that class.
\item An \emph{equivalent reformulation} is a target proved equivalent to the open problem. It is not an impossibility.
\item An \emph{unaffordable} route is a budget decision.
\item An \emph{unsuccessful attempt} claims nothing.
\end{itemize}

\subsection{Proved obstructions}
\begin{obstruction}[\cid{X.1}; elementary]\label{obs:support}
Integral sources do not bound prime support. Two examples show this:
\begin{itemize}[leftmargin=*]
\item The maps $\Z\to\Zp$, $1\mapsto p^2$, are nonzero at every prime and have nonunit image at every prime.
\item The two-term complex over $\prod_p\Zp$ with differential $(p^2)_p$ is perfect, yet has finite nonzero defects at infinitely many primes.
\end{itemize}
Hence source integrality, profinite perfectness, pointwise nonvanishing and pointwise finite primary groups do not, individually or together, imply finite support. By contrast, one fixed nonzero integer $C$ with $\Sha[p]\neq0\Rightarrow p\mid C$ would give finite support.
\end{obstruction}

\begin{obstruction}[\cid{X.2}; elementary]\label{obs:det}
Suppose the localization matrix of $\langle P_1,P_2\rangle/p$ is invertible. That says nothing about the quotient $\Sel_p/(E(\Q)/p)$: in an exact sequence $0\to M\to M\oplus Q\to Q\to0$, an isomorphism on $M$ carries no information about $Q$.
\end{obstruction}

\begin{researchrecord}[\cid{X.5}; written proof, in the stated class]\label{obs:mellin}
Proposition~\ref{prop:m1} has no realization inside the finite logarithmic class L3. Four facts are established:
\begin{itemize}[leftmargin=*]
\item Nonconstant polynomials in $\ell(z)$ have infinite cusp-translation span.
\item The ordinary algebraic infinitesimal Tate deformation does not give the required jet.
\item The half-ray maps to a null-homotopic loop in $E$.
\item No finite algebraic connection, iterated integral or finite exact correction in the class L3 supplies the required local identity.
\end{itemize}
The half-ray and full-ray moments satisfy $J_j=(1+(-1)^j)H_j$.
\end{researchrecord}

\begin{researchrecord}[\cid{X.6}; computer-assisted]\label{obs:jet}
A normalized, Atkin--Lehner-invariant second derivative $\mathcal E''(z,0)$ of real-analytic Eisenstein series satisfies
\[
\int_\gamma\alpha_f\,\mathcal E''(z,0)=\tfrac12L''(E,1)+\Delta ,
\]
with an explicit Fourier residual $\Delta$. Two single-vertex completions have certified residuals:
\begin{itemize}[leftmargin=*]
\item the midpoint maximal-compact completion: $0.000535614030\le\Delta_{\mathrm{mid}}\le0.000801512867$;
\item the Fricke-symmetrized endpoint completion: $0.000757821864\le\Delta_*\le0.000841639764$.
\end{itemize}
Both residuals are positive, so neither completion reproduces the weighted Mellin period. Arbitrary mixtures in the invariant span lie outside this class.
\end{researchrecord}

\begin{researchrecord}[\cid{X.7}; written proofs, each in its stated scope]\label{obs:midpoint}
Four further analyses of the midpoint residual each stop in their stated scope:
\begin{itemize}[leftmargin=*]
\item The $U_{4547}$-trace forces the translated residual periods to sum to zero, but it does not force the zeroth period to vanish.
\item The archived calculation excludes the specified finite combinations of ordinarily convergent Euler-product Dirichlet families and their derivatives by a Mellin-transform uniqueness argument with a nonzero continuous density. The precise representation and convergence class is part of that record; this summary excludes neither every conceivable Euler-product identity nor isolated special-value identities.
\item The residual matches the Berndt--K\"uhn weight-two series after differentiation, but it has no arithmetic Green-current identification.
\item The relative/holomorphic Deligne-regulator interpretation on compact $X_0(N)$ is obstructed: the curvature is nonzero and has no smooth extension across the cusp.
\end{itemize}
\end{researchrecord}

\begin{obstruction}[\cid{X.8}; written proof]\label{obs:curv}
Let $J$ be a global $C^1$ scalar function on $Y_0(N)$ such that $J$ and its first derivatives have polynomial-logarithmic growth at every cusp. Then for every holomorphic weight-two cusp form,
\[
\int_{Y_0(N)}\alpha_f\wedge dJ=0 .
\]
\end{obstruction}
\begin{proof}
$\alpha_f\wedge dJ=d(-J\alpha_f)$, and the cusp fluxes vanish by cusp decay. An independent unfolding of the Eisenstein family gives the same result.
\end{proof}
The class excludes singular, multivalued and nonexact kernels, and kernels with nonzero boundary flux.

\begin{obstruction}[\cid{X.10}; elementary]\label{obs:torus}
Let $k$ be a number field, $A/k$ an abelian variety, and $0\to\mathbf G_m^s\to\mathcal G\to A\to0$ a commutative extension. Then
\[
\Sha^1(k,\mathcal G)\xrightarrow{\ \sim\ }\Sha(k,A).
\]
The same holds for vector-group extensions and for quasi-trivial tori. For tori split by an extension of degree prime to $\ell$, the $\ell$-primary parts agree. In particular, the extension of $E$ by the split torus $\mathbf G_m^2$ attached to $P_1,P_2$ retains exactly the unknown group; it is also anti-affine \cite{brion}.
\end{obstruction}

\begin{obstruction}[\cid{X.11}; elementary]\label{obs:local}
Totally positive quadratic integers can approach $36$ while matching arbitrarily much prescribed finite local data. The first-derivative formulas tested in the records have vanishing inputs for this rank-two setting. The bare factorization $L(E/K,s)=L(E,s)L(E^K,s)$ supplies no arithmetic second-derivative comparison by itself. These are scoped tests, not exclusions of every future formula.
\end{obstruction}

\begin{obstruction}[\cid{X.13}; elementary]\label{obs:ssunit}
At a supersingular $p\ge5$, $a_p=0$. So $X^2+p$ has no unit root, and ordinary stabilization cannot be used; one has instead $\mathrm{pr}^2(\theta_{n+2})=-p\,\theta_n$. The signed objects of Theorem~\ref{thm:a3} replace the unit-root lane.
\end{obstruction}

The compression obstruction for the $2243$ certificate (\cid{X.15}) is Proposition~\ref{prop:compress}.

\subsection{Failure of the proposed trace-divisibility inference}
\begin{obstruction}[\cid{X.3}; counterexample]\label{obs:dm}
The unrestricted implication tested in the project is
\[
\alpha\in\mathfrak P,\ \mathfrak P\mid p\ \Longrightarrow\ \mathrm{Tr}(\alpha)\in p\Z .
\]
This is false: $4+i\in\Z[i]$ lies above $17$ but has trace $8$. This elementary counterexample alone is not a source-level refutation under an unstated hypothesis set. The separate Destrempes--Malinin audit is retained in the claim index with its source-specific scope. No exceptional-prime set is imported from that preprint here.
\end{obstruction}

\subsection{Equivalent reformulations}
\begin{researchrecord}[\cid{X.9}; written deduction applying published theorems]\label{obs:vis}
Varying-level visibility covers every relevant class for each fixed residual prime (Tamiozzo \cite{tamiozzo}, with Zhang's hypotheses \cite{zhang}). It gives no horizontal gain, for three reasons:
\begin{itemize}[leftmargin=*]
\item An auxiliary prime $\ell$ serves $p$ only if $p\mid D_\ell=(\ell+1)^2-a_\ell^2$, so a finite stock of auxiliary primes serves finitely many $p$.
\item Functorial descent to a fixed ambient variety that preserves the class up to a $p$-unit is obstructed.
\item For a fixed embedding, eventual fixed-ambient visibility on the ordinary tail is \emph{equivalent} to eventual ordinary vanishing.
\end{itemize}
\end{researchrecord}

\begin{researchrecord}[\cid{X.12}; elementary]\label{obs:stark}
Over $\Q(\sqrt{-23})$, where $\rank E(K)=3$, a fixed rank-three ambient base-change model within the finitely generated module class is equivalent to finite lane support. A bounded variant gives only $\#\Sha[p]\le p^{B-3}$.
\end{researchrecord}
Research records~\ref{thm:o1} and~\ref{thm:divfield} are equivalences of the same kind.

\subsection{Unsuccessful attempts and procedural stops}
\begin{itemize}[leftmargin=*]
\item \emph{Eisenstein congruences} (\cid{X.4}). Nonclassical weight-two methods reach only the Selmer subspace of extensions that factor through a chosen automorphic quotient, and nothing bounds the complement. Separately, $1\oplus\rho_E$ has Hodge--Tate weights $\{0,0,1\}$, which are irregular for classical $\GL(3)$ cohomology. The defect is one of coverage, not of arrow direction.
\item \emph{The prime-by-prime ladder} (\cid{X.14}). It was stopped at $13$. Climbing the ladder in order ($17$, $19$, $23$, \dots) moves only the conditional threshold ($2601\to3249\to4761\to\cdots$) and never reaches the infinite tail; a prime settled out of order does not move it, because the threshold advances only when the smallest remaining prime is excluded.
\item \emph{Other named constructions} (\cid{X.17}). Each was stopped without an impossibility theorem:
\begin{itemize}[leftmargin=*]
\item the cross-cusp kernel RS2 was archived by a stopping rule;
\item for the Heegner--Stark ideals $I_3$, no method of evaluation was found;
\item the witness $c_2(p)$ remains open at every prime;
\item a Saito--Kurokawa tangent inequality points in the wrong direction for this use;
\item the generalized Kato classes of the H9 program lack a comparison with the rational points;
\item one Castella--Hsieh branch misses the calibrated line at $p=5$.
\end{itemize}
\end{itemize}

\subsection{Unaffordable}
The current cost models estimate the $2243$ and $5843$ certificates (\cid{X.16}) at about $1.1\cdot10^5$ and $6.3\cdot10^5$ core-hours; these are model estimates, not necessary workloads. They are parked by decision, and the $2243$ evaluator's restart defect is to be fixed first. This is not a mathematical stop.
